\subsection{Flow-Matching Objective}
Given a ground-truth action chunk $a$ and base noise $z \sim \mathcal{N}(0,I)$, we construct an interpolated state
\[
x_\tau = (1-\tau)z + \tau a, \quad \tau \sim \mathcal{U}(0,1),
\]
with corresponding target velocity
\[
u = a - z .
\]
The action expert $v_\phi$ is trained via the flow-matching objective
\[
\mathcal{L}_{\text{flow}} =
\mathbb{E}\!\left[
\big\| v_\phi(x_\tau,\tau; o, l) - (a - z) \big\|_2^2
\right],
\]
which encourages the model to learn a continuous vector field transporting noise to valid actions conditioned on the observation and task.

At inference time, actions are generated by integrating the learned dynamics
\[
\frac{dx_\tau}{d\tau} = v_\phi(x_\tau,\tau; o, l)
\]
from $\tau=0$ to $\tau=1$ using a fixed-step explicit ODE solver.


\subsection{Training Data} 
The full dataset mixture used for training is summarized in Table~\ref{tab:data_mix}. The reported weights account for each dataset’s relative size and therefore reflect the true fraction of samples drawn from each source within a training batch. We assign a higher proportion to the DROID dataset~\cite{khazatsky2025droidlargescaleinthewildrobot} due to its broad task coverage and rich environmental diversity.

\begin{table}[th]
    \centering
    \caption{Training data mixture. Percentages denote the true fraction of samples from each dataset within each training batch.}
    
    \begin{tabular}{lr}
        \toprule
        \multicolumn{2}{c}{\textbf{Training Dataset Mixture}}\\
        \midrule
        DROID \cite{khazatsky2025droidlargescaleinthewildrobot} & 85.26\% \\
        Fractal \cite{brohan2023rt1roboticstransformerrealworld} & 5.86\% \\
        Bridge \cite{walke2024bridgedatav2datasetrobot} & 3.39\% \\
        BC-Z \cite{pmlr-v164-jang22a} & 0.47\% \\
        Taco Play \cite{mees2023groundinglanguagevisualaffordances} & 0.73\% \\
        Jaco Play \cite{dass2023jacoplay} & 0.12\% \\
        Furniture Bench Dataset \cite{heo2023furniturebenchreproduciblerealworldbenchmark} & 0.31\% \\
        UTAustin Mutex \cite{shah2023mutex} & 0.56\% \\
        Berkeley Fanuc Manipulation \cite{zhu2023fanuc} & 0.19\% \\
        FMB Dataset \cite{luo2024fmbfunctionalmanipulationbenchmark} & 0.09\% \\
        Berkeley Autolab UR5 \cite{BerkeleyUR5Website} & 0.15\% \\
        Austin Buds Dataset \cite{zhu2022bottomupskilldiscoveryunsegmented} & 0.05\% \\
        Austin Sailor Dataset \cite{nasiriany2022sailor} & 0.53\% \\
        Austin Sirius Dataset \cite{liu2023robotlearningjobhumanintheloop} & 0.44\% \\
        Viola \cite{zhu2023violaimitationlearningvisionbased} & 0.12\% \\
        MolmoAct \cite{lee2025molmoactactionreasoningmodels} & 1.73\% \\
        \bottomrule
    \end{tabular}
    \label{tab:data_mix}
\end{table}


During training, we filter idle segments by removing action sequences below a small motion threshold for consecutive steps and discard trajectories lacking task instructions.

\subsection{Compute and Infrastructure Details}
\label{app:compute}

\noindent\textbf{Hero-run compute.}
The final reported checkpoints for all models, including \textsc{LAP-3B}, \textsc{LAP-3B} with VQA co-training, and all replicated baselines, were trained using an identical large-scale pre-training setup. Each hero run was trained for approximately 50 hours on a TPU v6e-64 slice, corresponding to the checkpoints used for all main experiments. Across all reported models, the total hero-run compute amounts to approximately \textbf{50 TPU v6e-64 hours per model}.

\medskip
\noindent\textbf{Total pre-training compute.}
Across the full set of pre-training experiments—including \textsc{LAP-3B} and all replicated baselines, but excluding early prototyping, debugging runs, and downstream fine-tuning experiments—we conducted approximately 200 large-scale training runs. Collectively, these experiments consumed over \textbf{4,000 TPU v6e-64 hours}.

\medskip
\noindent\textbf{Inference compute.}
All real-robot evaluations are executed in real time using a single commodity GPU (an NVIDIA RTX 4090), without any additional model parallelism, specialized hardware, or acceleration techniques. End-to-end real-robot evaluation for the reported experiments required approximately \textbf{24 hours of total human-supervised runtime}.



\section{Evaluation Protocol}
\label{app:evaluation_details}

Each task is evaluated over 20 independent trials with randomized object placement within workspace bounds. A trial is terminated upon task success, time-out, or a safety-triggered abort, such as when the robot collides with the table, reaches joint limits, or exhibits unsafe behavior (e.g., excessively rapid or unstable motions).

For zero-shot experiments, we report binary task success along with 95\% finite-sample valid confidence intervals for binomial success rates, which control the Type-I error (miscoverage) probability \cite{vincent2024generalizablebehaviorcloningpolicy}.

For fine-tuning experiments, we use staged success scores to capture partial task completion:

\noindent$\bullet$ \emph{Hang Tape on Rack:}
\begin{itemize}
\item 0.25: reach the tape
\item 0.5: grasp the tape
\item 0.75: reach the rack while holding the tape
\item 1.0: successfully hang the tape
\end{itemize}

\noindent$\bullet$ \emph{Fold Towel and Place in Basket:}
\begin{itemize}
\item 0.5: correctly folded towel
\item 0.75: placed in basket but unfolded
\item 1.0: full task success
\end{itemize}




\begin{figure*}[t]
    \centering
    \includegraphics[width=\linewidth]{figures/appendix/fig9_rollout_visualization.pdf}
    \caption{\textbf{Zero-shot evaluation rollout visualizations.} Example observations from the camera views used for policy evaluation across different robot embodiments and tasks, shown from left to right.}
    \label{fig:zeroshot-rollout-visualization}
\end{figure*}

